\section{Phân Tích Thay Đổi Trong App Repo \texttt{tzin1401/yas}}

App repo không chỉ là source code gốc của YAS. Trong đồ án này, app repo đã được
thay đổi để trở thành điểm xuất phát của CD: phát hiện module thay đổi, chạy CI
gates, build Docker image, push Docker Hub, rồi cập nhật CD repo bằng GitOps.
Các nội dung dưới đây cần được đưa vào báo cáo vì chúng giải thích vì sao CD repo
có thể deploy đúng image và đúng service.

\subsection{Vai Trò Của App Repo Trong Kiến Trúc Split-Repo}

App repo là nơi tạo artifact deployable, nhưng không phải nơi lưu trạng thái mong
muốn cuối cùng của cluster. Contract giữa app repo và CD repo gồm ba dữ liệu:

\begin{table}[H]
\centering
\caption{Contract app repo xuất sang CD flow}
\begin{tabular}{|p{3.2cm}|p{5.2cm}|p{5.4cm}|}
\hline
\textbf{Dữ liệu} & \textbf{Nguồn trong app repo} & \textbf{Cách CD repo sử dụng} \\
\hline
Service catalog & \path{services.yaml} & Jenkins biết service nào deployable, image name, Dockerfile và dependency expansion. \\
\hline
Image tag & Branch/main/tag context trong Jenkinsfile & CD repo overlay nhận tag commit SHA hoặc release tag immutable. \\
\hline
CI verdict & Gitleaks, test, coverage, SonarQube, build result & Chỉ image đã qua quality gate mới được dùng làm nguồn deploy/promote. \\
\hline
\end{tabular}
\end{table}

\begin{figure}[H]
\centering
\reportimg{img/jenkins/app_repo_pipeline_overview.png}
\caption{Giao diện tổng quan các stage trong pipeline CI của App Repo}
\end{figure}

\subsection{Jenkinsfile Chính: CI Gate Và CD Trigger}

Jenkinsfile trên app repo \texttt{origin/main} dùng agent label \texttt{yas-build-worker}
cho các bước nặng như test, build, scan và Docker. Pipeline chính có các nhóm stage:

\begin{itemize}
    \item Checkout, fetch \texttt{origin/main}, chmod Maven wrapper và kiểm tra tool.
    \item Detect changed modules bằng \texttt{ci/detect-changed-modules.sh}.
    \item Gitleaks, test, JaCoCo report, coverage gate và SonarQube.
    \item Docker Hub build/push image cho các service deployable.
    \item GitOps update: clone \texttt{emanhthangngot/yas-cd}, chỉnh overlay, commit và push vào \texttt{main}.
\end{itemize}

Một điểm phải ghi rõ: stage Snyk hiện đang bị tắt tạm thời bằng điều kiện
\texttt{when \{ expression \{ false \} \}} để unblock CD validation. Vì vậy báo cáo
không claim Snyk đang chạy ở bản app repo hiện tại; chỉ ghi đây là CI gate kế thừa
từ Đồ án 1 và cần bật lại sau khi CD ổn định.

\subsection{Chiến Lược Tag Image}

\begin{table}[H]
\centering
\caption{Chiến lược image tag trong app repo}
\begin{tabular}{|p{3.2cm}|p{5.2cm}|p{5.2cm}|}
\hline
\textbf{Nguồn trigger} & \textbf{Image tag được push} & \textbf{GitOps target} \\
\hline
Feature branch & Short commit SHA 12 ký tự & Không update GitOps trực tiếp; developer preview dùng job riêng \\
\hline
\texttt{main} & Short commit SHA, \texttt{main}, \texttt{latest} & Update \texttt{overlays/dev} bằng short commit SHA \\
\hline
\texttt{vX.Y.Z} tag & Retag image hiện có sang \texttt{vX.Y.Z} & Update \texttt{overlays/staging} bằng immutable release tag \\
\hline
\end{tabular}
\end{table}

Việc \texttt{dev} dùng commit SHA thay vì \texttt{main/latest} là thay đổi quan trọng:
mỗi lần main build thành công đều tạo Deployment template diff thật, ArgoCD có lý do
rollout pod mới, đồng thời rollback có thể truy ngược về tag bất biến.

\subsection{Release Promote Từ Dev Overlay}

Trong giai đoạn đầu, release tag cố promote mọi service từ cùng một tagged commit
SHA. Cách này lỗi khi \texttt{dev} đang là trạng thái incremental: một số service
đang ở SHA mới, một số service vẫn ở tag cũ. App repo đã sửa Jenkinsfile để đọc
tag nguồn từng service từ \path{overlays/dev/kustomization.yaml} trong CD repo,
sau đó dùng \texttt{docker buildx imagetools create} retag từng image sang release
\texttt{vX.Y.Z}. Nhờ đó release \texttt{v0.1.3} promote thành công qua Jenkins tag pipeline.

\subsection{Service Catalog Trong App Repo}

File \path{services.yaml} trong app repo là contract cho Jenkins:

\begin{itemize}
    \item Xác định service nào deployable, Dockerfile nằm ở đâu, imageName là gì.
    \item Ghi demo role: teacher-required, runtime-dependency, optional, run-once.
    \item Ghi dependencies để CI có thể mở rộng phạm vi build/test khi dependency thay đổi.
\end{itemize}

{\raggedright
Các service deployable chính là \texttt{cart}, \texttt{customer}, \texttt{inventory},
\texttt{location}, \texttt{media}, \texttt{order}, \texttt{payment},
\texttt{payment}-\allowbreak\texttt{paypal}, \texttt{product}, \texttt{search}, \texttt{tax},
\texttt{backoffice}-\allowbreak\texttt{bff}, \texttt{storefront}-\allowbreak\texttt{bff}, \texttt{backoffice}-\allowbreak\texttt{ui},
\texttt{storefront}-\allowbreak\texttt{ui}. Các service \texttt{promotion}, \texttt{rating},
\texttt{recommendation}, \texttt{webhook} được giữ trong source nhưng không thuộc
runtime CQ mặc định. \texttt{sampledata} được xử lý theo cơ chế run-once, không
chạy liên tục sau khi seed dữ liệu.\par}

\subsection{Changed-Module Detection}

Script \texttt{ci/detect-changed-modules.sh} giúp pipeline không rebuild toàn bộ
monorepo khi chỉ đổi một module. Logic chính:

\begin{itemize}
    \item Nếu chỉ đổi docs/spec/GitOps metadata thì trả về \path{__skip_full_ci__}.
    \item Nếu đổi Jenkinsfile, root \texttt{pom.xml}, \texttt{services.yaml}, Docker/K8s scripts thì chạy full path.
    \item Nếu đổi một module, chọn module đó và các service phụ thuộc dựa trên \texttt{services.yaml}.
\end{itemize}

Phần này đáng đưa vào báo cáo vì nó giải thích vì sao CD có thể scale trên monorepo
YAS lớn mà không tốn thời gian build mọi service trong mỗi commit.

\subsection{Build Và Runtime Fixes}

App repo cũng có các thay đổi hỗ trợ CD chạy ổn định:

\begin{itemize}
    \item Dockerfile của backend được chỉnh để copy đúng Spring Boot jar.
    \item \texttt{payment-paypal} được đóng gói đúng executable jar/thin jar theo nhu cầu runtime.
    \item Search integration tests được cô lập để tránh xung đột Testcontainers.
    \item Payment, product, tax, location, rating, webhook bổ sung test để giữ coverage gate.
\end{itemize}

Những thay đổi này không phải phần GitOps, nhưng là điều kiện để image sinh từ app
repo có thể boot được trên K3s.

\subsection{Observability Instrumentation Trong App Repo}

App repo đã bổ sung nền observability ở tầng ứng dụng:

\begin{itemize}
    \item Root \texttt{pom.xml} thêm \texttt{spring-boot-starter-actuator},
    \texttt{micrometer-registry-prometheus} và \texttt{spring-boot-starter-opentelemetry}.
    \item Nhiều service expose \path{/actuator/prometheus} qua
    \path{management.endpoints.web.exposure.include=prometheus}.
    \item SecurityConfig cho phép Prometheus scrape \texttt{/actuator/prometheus}
    và health probes mà không cần login.
    \item App repo có \texttt{docker-compose.o11y.yml}, Prometheus config, Grafana
    datasource/dashboard, Loki, Tempo và OpenTelemetry Collector cho local observability.
\end{itemize}

Vì vậy chương Observability của báo cáo phải phân biệt hai phần: app repo đã có
instrumentation; CD/K8s cần runtime evidence Prometheus/Grafana để chứng minh scrape
và dashboard hoạt động trên cluster.
